
\documentclass[10pt, twoside]{article}
\usepackage{amsfonts,amssymb,latexsym,xspace}
\usepackage{amsmath,enumerate}
\numberwithin{equation}{section}

%\usepackage{showkeys}

\hfuzz 15pt
\newtheorem{thm}{Teorema}[section]
\newtheorem{lem}{Lemma}[section]
\newtheorem{cor}{Corollario}[section]
\newtheorem{sublem}{Sub-Lemma}[section]
\newtheorem{prop}{Proposizione}[section]
\newtheorem{defin}{Definizione}
\newtheorem{cond}{Condizione}
\newtheorem{exmp}{Esempio}
\newtheorem{rem}{Remark}[section]

\newenvironment{proof}{{\sc Dimostrazione\ }}{\hfill $\Box$\vskip.6pt}
\arraycolsep=.5\arraycolsep

\newcommand\Cal[1]{{\mathcal #1}}

\newcommand\B{{\Cal B}}
\newcommand\Co{{\Cal C}}
\newcommand\D{{\Cal D}}
\newcommand\F{{\Cal F}}
\newcommand\Lp{{\Cal L}}
\newcommand\M{{[0,1]}}
\newcommand\T{{\Cal T}}
\newcommand\U{{\Cal U}}
\newcommand\V{{\Cal V}}

\newcommand\Ve{{\mathbb V}}
\newcommand\N{\mathbb N}
\newcommand\R{\mathbb R}
\newcommand\Z{\mathbb Z}
\newcommand\Q{\mathbb Q}
\newcommand\C{{\mathbb C}}
\newcommand\E{{\mathbb E}}
\newcommand\Pe{{\mathbb P}}
\newcommand\cp{{\Cal P}}

\newcommand\BV{{\text{BV}}}
\newcommand\vf{\varphi}
\newcommand\nn{\nonumber}
\newcommand\ve{\varepsilon}

\pagestyle{myheadings}
\markboth{\centerline{\sc Carlangelo Liverani}\hskip-12pt}
{\centerline{\sc Numeri razionali e reali}\hskip-12pt}

\begin{document}
\title{\bfseries Numeri razionali e reali}
\author{{\sc Carlangelo Liverani}\\
Dipartimento di Matematica\\
Universit\`{a} di Roma (Tor Vergata)\\
Via della Ricerca Scientifica, 00133 Roma, Italy\\
{\tt liverani@mat.uniroma2.it}}
\date{Rome, Ottobre 23, 2001}
\maketitle


\section{\sc \normalsize Numeri Razionali}

\begin{defin}
Dati $p,q\in\N$ diciamo che \emph{$p$ \`e divisibile per $q$} e
scriviamo $q|p$ se esiste $c\in \N$ tale che $p=cq$.
\end{defin}

\begin{lem}\label{lem:divisibility}
Se $p\in\N$ \`e primo e se , dati due numeri $a,b\in\N$, $p|ab$ allora
o $p|a$ oppure $p|b$.
\end{lem}
\begin{proof}
Chiaramente \`e sempre possibile trovare $m,n,a_1,b_1\in\N$,
$a_1,b_1<p$, tali che 
\begin{equation}\label{eq:reduction}
\begin{split}
a&=pn+a_1\\
b&=pm+b_1.
\end{split}
\end{equation}
Il Lemma afferma che o $a_1$ oppure $b_1$ devono essere
nulli, cio\`e $a_1b_1=0$. Per dimostrare che deve essere cos\`\i\ 
ragioniamo per assurdo. Supponiamo che $a_1b_1\neq 0$,
allora (\ref{eq:reduction}) implica $ab= p^2nm+pm aa_1+pnb_1+a_1b_1$,
da cui segue che $p|a_1b_1$. Il Lemma \`e dunque equivalente a dire
che per ogni $p$ primo non esistono $a,b<p$ tali che $p|ab$.

Se $p=2$ il risultato \`e ovvio, dunque il Lemma \`e vero per alcuni
numeri primi. Ragioniamo nuovamente per assurdo e supponiamo che esista qualche
numero primo per cui il Lemma \`e falso. Sia $p'$ il pi\`u piccolo di
tali numeri primi. Dunque esisteranno $a,b\in\N$, $a,b<p'$ tali che
$p'|ab$. Ora sia $c\in \N$ tale che
$ab=cp'$, chiaramente $c<p'$. Dunque $c|ab$ e qualunque numero primo
che divida $c$ divider\`a anche $ab$. Ne segue che o $a$ oppure $b$
\`e divisibile per tale numero primo (poich\`e esso \`e
necessariamente pi\`u piccolo di $p'P$ e dunque \`e, per costruzione,
uno dei numeri primi per cui il Lemma \`e noto essere vero). Otterremo
perci\`o numeri $a_1,b_1, c_1$ tali che $a_1b_1=c_1p'$.
Continuando in questa maniera possiamo dividere per un divisore primo
di $c_1$ ottenendo $a_2,b_2,c_2$ per cui $a_2b_2=c_2 p'$ e cosi via
fino a che otteniamo $a_nb_n=p'$ (cio\`e $c_n=1$). Ma questa ultima
uguaglianza \`e impossibile poich\`e contraddice l'ipotesi che $p'$ \`e primo.
\end{proof}

\begin{lem}
Ogni numero intero $a\in\N$ ha una unica espressione della forma
$a=p_1^{k_1}\cdots p_n^{k_n}$ dove  $\{p_i\}$ sono numeri primi e
$\{k_i\}$ sono interi. 
\end{lem}
\begin{proof}
Supponiamo che esistano due tali decomposizioni $a=p_1^{k_1}\cdots
p_n^{k_n}$ e $a=q_1^{j_1}\cdots q_n^{j_m}$. Supponiamo inoltre che
esista $q_i\not\in\{p_1,\dots,p_n\}$, allora $q_i|p_1^{k_1}\cdots
p_n^{k_n}$ ma se applichiamo il Lemma \ref{lem:divisibility}
ripetutamente otteniamo che $q_i$ deve dividere uno dei $\{p_i\}$ cosa
assurda perch\`e essi sono primi. Dunque $\{q_i\}=\{p_i\}$. La sola
possibilit\`a che rimane \`e che $\{k_i\}\neq \{j_i\}$ ma dividendo
opportunamente ci si riconduce immediatamente al caso precedente.
\end{proof}
\begin{lem}
Ogni numero razionale ha una unica rappresentazione del tipo $\frac
pq$ con $p$ e $q$ primi fra loro.
\end{lem}
\begin{proof}
Supponiamo 
\[
\frac{p}{q}=\frac{p'}{q'}
\]
con $q<q'$. Sia $c$ un numero primo tale che $c|q'$. Allora
$qp'=pq'$ dunque $c|qp'$ e, per il Lemma
\ref{lem:divisibility} e ricordando che $q'$ e $p'$ sono primi tra
loro, segue $c|q$. Dunque $q=cq_1$, $q'=q_1'c$ e quindi $q_1p'=pq'_1$,
continuando a dividere per numeri primi si otterr\`a $p'=p$ da cui
seque $q=q'$.
\end{proof}
\section{\sc \normalsize Algoritmo Euclideo e frazioni continue}

Il problema aperto dalla sezione predente \`e quello di ridurre una
frazione alla sua forma standard, ovvero \emph{trovare il massimo
comun divisore tra due numeri dati}. Questo problema ha una elegante
soluzione conosciuta come \emph{algoritmo Euclideo}.

Siano $p,q\in\N$, $p>q$, allora esiste $a_1\in\N$ tale che
\begin{equation}\label{eq:rapp}
p=a_1q+p_1,\, \quad p_1<q
\end{equation}
e
\[
\frac pq= a_1+\frac 1{\frac {\textstyle q}{\textstyle p_1}}
\]
se $q$ e $p_1$ hanno un divisore comune (chiamiamolo $b$) allora anche
$b|p$ (da (\ref{eq:rapp})) e quindi $b$ \`e divisore comune di $p,
q$. D'altro canto se $b|p$ e $b|q$ allora $b|p_1$ (questa \`e ancora
una volta conseguenza di (\ref{eq:rapp})). Dunque $p,q$ hanno
gli stessi divisori comuni di $q,p_1$ e quindi lo stesso massimo comun
divisore. Poich\`e 
\[
q=a_2p_1+q_1,\, \quad q_1<p_1
\]
possiamo scrivere
\[
\frac pq=a_1+\frac 1{a_2+\frac 1{\frac{\textstyle p_1}{\textstyle q_1}}}
\]
dove $q_1,p_1$ hanno gli stessi divisori comuni di $q,p_1$ e quindi
anche di $p,q$. Questa
procedura pu\`o essere iterata fino a che non si ottiene un resto
nullo. Per esempio se $q_{n-1}=a_{2n}p_n$, cio\`e $q_n=0$, allora
$q_{n-1}$ e $p_n$ sono entrambi divisibili per $p_n$ e quindi $p_n$
(l'ultimo resto non nullo) \`e il massimo comun divisore.

Per di pi\`u abbiamo ottenuto un'altra maniera univoca di scrivere una
frazione chiamata \emph{frazione continua}:
\[
\frac pq= a_1+\frac 1{\textstyle a_2+{\ \atop\ddots {\
\atop{\textstyle +\frac {\textstyle 1}{\textstyle a_{2n}}}}}} 
\]

\begin{exmp}
Si trovi il massimo comun divisore tra $p=855$ e $q=266$.
\[
\begin{split}
855=&3\cdot 266+57\\
266=&4\cdot 57+38\\
57=&1\cdot 38+19\\
38=&2\cdot 19.
\end{split}
\]
Dunque il massimo comun divisore \`e 19 e si ha la seguente
rappresentazione in frazione continua.
\[
\frac{855}{266}=3+\frac {\textstyle 1}{\textstyle 4+\frac {\textstyle
1}{\textstyle 1+\frac{\textstyle 1}{\textstyle 2}}}
\]
\end{exmp}

\section{\sc \normalsize I Numeri Razionali non bastano}

Allo stesso modo in cui i numeri razionali devono essere introdotti
per attribuire un 
risultato ad ogni divisione, i numeri irrazionali emergono
naturalmente nello studio della operazione inversa della potenza:
l'estrazione di radice. Il caso pi\`u semplice che si puo considerare
\`e il seguente: esiste un numero $a$ tale che $a^2=2$? (Stiamo
parlando di $\sqrt 2$.)

\begin{lem}
Non esiste alcun numero razionale tale che il suo quadrato sia 2.
\end{lem}
\begin{proof}
Supponiamo che esistano $p,q\in \N$, primi tra loro, tali che 
\[
\left(\frac p q\right)^2=2.
\]
Allora $p^2=2q^2$, dunque per il Lemma \ref{lem:divisibility}
$2|p$. Dunque, $p=2p_1$ e quindi $4p_1^2=2q^2$ che implica
$q^2=2p_1^2$. Ma allora, sempre per il Lemma \ref{lem:divisibility},
$2|q$ cosa che contraddice l'ipotesi che $p,q$ siano primi tra loro. 
\end{proof}

La cosa \`e seccante ma, come premio di consolazione, esistono numeri
razionali molto vicini a quello voluto.
\begin{lem}\label{lem:approx2}
Per ogni $n\in\N$ esiste un numero razionale $a\in\Q$ tale che
\[
|a^2-2|\leq \frac 1{n}.
\]
\end{lem}
\begin{proof}
Sia $a_0=1$ e $a_1=2$, chiaramente $a_0^2<2$ e $a_1^2>2$. Consideriamo
il punto intermedio $a_2=\frac{a_0+a_1}2=\frac 32$, risulta che
$a_2^2>2$. Allora $2\in[a_0^2,a_2^2]$. Possiamo nuovamente considerare
il punto intermedio $a_3=\frac{a_0+a_2}2$ e verificare se $a_3^2$ \`e
maggiore o minore di due (accade che $a_3^2<2$). Possiamo dunque dire
che $2$ \`e contenuto nell'intervallo pi\`u stretto
$[a_3^2,a_2^2]$. Si pu\`o continuare nello stesso modo quante volte si
vuole (diciamo $m$ volte) ottenedo cosi $2\in[a_k^2,a_j^2]$ con
$|a_k^2-a_j^2|\leq 2^{-m+2}$. Basta dunque che $2^{m-2}>n$ per
dimostrare il Lemma.\footnote{Se per esempio vogliamo conoscere un
numero razionale $a$ tale che $|a^2-2|\leq 10^{-3}$ allora, poich\`e
$2^{10}=1024$, baster\`a iterare la procedura descritta 12 volte.}
\end{proof}

Una cosa interessante della dimostrazione del Lemma precedente \`e che
si tratta di una dimostrazione costruttiva, cio\`e viene specificata
una procedura esplicita per costruire i numeri richiesti. Tuttavia la
procedura specificata non \`e n\`e l'unica n\`e la pi\`u
efficiente. Nuovamente ci troviamo davanti ad una
ambiguit\`a. Ciononostante, se numeri diversi soddisfano alle ipotesi
del Lemma \ref{lem:approx2} allora devono essere molto vicini tra
loro.
\begin{lem}
Se $a,b \in \Q$ sono tali che $|a^2-2|\leq \frac 1n$ e $|b^2-2|\leq
\frac 1n$, per qualche $n\in\N$, allora
\[
|a-b|\leq \frac 1n.
\]
\end{lem}
\begin{proof}
Si noti che
\[
|a^2-b^2|\leq |a^2-2|+|b^2-2|\leq \frac 2n.
\]
Per di pi\'u $a^2-2\geq -1/n\geq -1$ implica $a^2\geq 1$, cio\`e
$a,b>1$. Dunque $|a^2-b^2|=(a+b)|a-b|\geq 2|a-b|$. Collezionando le
precedenti disuguaglianze abbiamo $|a-b|\leq \frac 1n$.
\end{proof}

Per quanto detto prima \`e interessante vedere una dimostrazione
alternativa del Lemma  \ref{lem:approx2}. Per costruire un'altra
successione di numeri con la propriet\`a desiderata applichiamo
formalmente l' algoritmo Euclideo al misterioso numero $\sqrt 2$.
\begin{equation}\label{eq:magic}
\sqrt 2=\frac{\sqrt 2}1= 1+\frac{\textstyle 1}{\textstyle
\frac{\textstyle 1}{\textstyle \sqrt 2-1}}=1+\frac{\textstyle 1}{\textstyle
\frac{\textstyle \sqrt 2+1}{\textstyle  2-1}}= 1+\frac{\textstyle
1}{\textstyle 1+\sqrt 2}.
\end{equation}
Iterando la equazione (\ref{eq:magic}) si ottiene
\[
\sqrt 2=1+\frac{\textstyle 1}{\textstyle 2+\frac{\textstyle
1}{\textstyle 1+\sqrt 2}}
\]
e, usando nuovamente (\ref{eq:magic}),
\[
\sqrt 2=1+\frac{\textstyle 1}{\textstyle 2+\frac{\textstyle
1}{\textstyle 2+\frac{\textstyle
1}{\textstyle 1+\sqrt 2}}}
\]

Le formule precedenti suggeriscono un nuovo modo di costruire una
sequenza di numeri il cui quadrato si avvicini sempre pi\`u a $2$. Nel
seguito renderemo rigorosa e chiarificheremo tale intuizione.

La formula (\ref{eq:magic}) suggerisce di definite la funzione
$f:\Q\to\Q$
\begin{equation}\label{eq:frazdef}
f(x):=1+\frac {1}{1+x}.
\end{equation}
Studiamo le propriet\`a della funzione $f$.

Prima di tutto notiamo che 
\[
f(x)^2=1+\frac{3+2x}{1+2x+x^2}.
\]
Da questa formula segue immediatamente che se $x^2>2$ allora $f(x)^2<2$ e se
$x^2<2$ allora $f(x)^2>2$. Inoltre
\[
\begin{split}
f(1)=&\frac 32<2\\
f(2)=&\frac 43>1.
\end{split}
\]
Infine la funzione $f$ inverte l'ordine, cio\`e se $0<x<y$ allora
\[
f(x)=1+\frac 1{1+x}>1+\frac 1{1+y}=f(y).
\]
Tutto cio significa che $f([1,2])\subset [1,2]$ e che se noi prendiamo
un qualunque numero $a\in[1,2]$ allora $f(a)$, $f^2(a)=f(f(a))$,
$f^3(a)=f(f(f(a)))$ sono tutti numeri nell'intervallo $[1,2]$. Per di
pi\`u continuano a saltare a destra ed a sinistra di $\sqrt 2$. La
domanda che rimane \`e se questi numeri si avvicinano a $\sqrt 2$ oppure no.

Siano $x,y\in[1,2]$ allora
\[
|f(x)-f(y)|=\left|\frac
1{1+x}-\frac1{1+y}\right|=\frac{|x-y|}{(1+x)(1+y)}\leq \frac{|x-y|}4.
\]
Dunque $f^1([1,2])$ \`e un intervallo di lunghezza $1/4$; $f^2([1,2])$
\`e un intervallo di lunghezza  $4^{-2}$ e $f^n([1,2])$
\`e un intervallo di lunghezza  $4^{-n}$. D'altro canto se $a<b$ sono
gli estremi dell'intervallo $f^n([1,2])$ allora deve essere $a^2<2$ e
$b^2>2$. Questo dimostra che scegliendo, per esempio, $a_0=1$,
$a_{n+1}=f(a_n)$ si ottiene un'altra successione del tipo richiesto
dal Lemma \ref{lem:approx2}, solo questa volta la rapidit\`a con cui
si raggiunge una data precisione \`e doppia.
Per esempio
\[
\begin{split}
a_1&=f(a_0)=f(1)=1+\frac 12=\frac 32\;;\quad
a_2=f(a_1)=f(3/2)=1+\frac{\textstyle 1}{\textstyle 2+\frac{\textstyle
1}{\textstyle 2}}=\frac 75\\
a_3&=f(a_2)=f(7/5)=1+\frac{\textstyle 1}{\textstyle 2+\frac{\textstyle
1}{\textstyle 2+\frac{\textstyle
1}{\textstyle 2}}}=\frac {17}{12}
\end{split}
\]
e cosi via.
Concludendo tutta questa discussione, siamo arrivati alla comprensione
del fatto che $\sqrt 2$ \`e un \emph{numero} ben definito, infatti si
pu\`o parlare senza ambiguit\`a della sua $n$-esima cifra
decimale. Tuttavia risulta che $\sqrt 2$ \`e definibile
solo attraverso una procedura di approssimazione che ci permetta di calcolarne
\emph{tante cifre quante vogliamo} (almeno in linea di
principio). Queste procedure di 
approssimazione non sono uniche (ne abbiamo viste due differenti ma ne
esistono infinite altre) ma sono tutte equivalenti, nel senso che
stabilita la precisione richiesta danno tutti numeri che differiscono
tra di loro per un errore inferiore alla quantit\`a specificata.
Queste procedure di approssimazione sono formalizzate nel concetto di
\emph{limite} di cui discuteremo tra poco. L'insieme dei \emph{numeri}
definibili attraverso tali procedure \`e detto \emph{insieme dei
numeri reali} (in simboli $\R$) e i numeri reali ma non razionali
($\{x\in\R\;|\; x\not\in \Q\}$) sono detti \emph{numeri irraziomali}.

\section{\sc \normalsize Numeri Irrazionali}

Nella sezione precedente abbiamo visto che $\sqrt 2$ pu\`o essere
\emph{approssimato con precisone arbitraria}. Questo concetto \`e
formalizzato in maniera precisa nella seguente definizione di limite.
\begin{defin}
Data una successione $\{a_n\}_{n\in\N}$, si dice che il limite di
$\{a_n\}_{n\in\N}$ \`e $b$, e si scrive $\lim_{n\to\infty}a_n=b$, se
per ogni $\varepsilon>0$ esiste $\bar n\in N$ tale che
\[
|a_n-b|\leq \varepsilon\quad \text{per ogni }n\geq \bar n.
\]
\end{defin}

Per esempio \`e facile verificare che $\lim_{n\to\infty}\frac
1n=0$. La cosa seccante della definizione precedente \`e che richiede
l'esistenza del numero $b$ e quindi non si pu\`o applicare alla nostra
precedente discussione, sebbene da essa siamo convinti che $\sqrt 2$
\`e un numero ben definito. Come abbiamo gi\`a detto questa
convinzione deriva dal fatto che possiamo calcolarne quante cifre
vogliamo. Anche questa propriet\`a pu\`o essere formalizzata
rigorosamente nel concetto di \emph{successioni di Cauchy}.
\begin{defin}
Si dice che una successione $\{a_n\}_{n\in\N}$ \`e di Cauchy se
per ogni $\varepsilon>0$ esiste $\bar n\in N$ tale che
\[
|a_n-a_m|\leq \varepsilon\quad \text{per ogni }n,m\geq \bar n.
\]
\end{defin}
Chiaramente le successioni che abbiamo trovato in relazione a $\sqrt
2$ sono Cauchy.

\`E facile verificare che se una successione di numeri razionali ha
limite allora \`e di Cauchy mentre, come abbiamo visto, il contrario
non \`e necessariamente vero. 

A questo punto possiamo aggiungere tutte le radici quadrate ai numeri razionali
ed ottenere un insieme di numeri in cui si pu\`o sempre fare
l'operazione di radice quadrata. Ma che dire delle radici cubiche? Il
processo di allargamento dei numeri sembra non avere fine e
apparentemente ci servono sempre pi\`u numeri man mano che
consideriamo nuove operazioni, la situazione \`e alquanto
intollerabile. D'altro canto la nostra discussione di $\sqrt 2$ ci ha
insegnato che si puo aggiungere un numdero solo se questo \`e
conoscibile con precisione arbitraria, altrimenti non si sa di che si
sta parlando.  Ci\`o significa che la soluzione del problema
dell'allargamento dei numeri deve risiedere nella richiesta che
\emph{ogni successione di Cauchy abbia limite}. Il modo ovvio di fare
ci\`o \`e di dire che i numeri sono appunto tali successioni di
Cauchy. Ovviamente diverse successioni di Cauchy possono dare luogo
allo stesso numero (lo abbiamo visto per $\sqrt 2$) ma questo non ci
pu\`o spaventare pi\`u di tanto visto che \`e una situazione che si
\`e gi\`a presentata con i numneri razionali rappresentati come
frazioni. Tuttavia questa ambiguit\`a per i numeri razionali ha dato
luogo a parecchie riflessioni per convincersi che era
innoqua. Non pu\`o dunque sorprendere che per rendere precisa l'idea di
completamento di cui sopra siano necessarie parecchie riflessioni che
tuttavia esulano dallo scopo del presente corso in cui ci
accontenteremo della precedente idea intuiva di \emph{numero reale}.

Si noti comunque che le varie propriet\`a dei limiti (e.g. il limite
della somma \`e la somma dei limiti, lo stesso per il prodotto,
sottrazione, etc.) ci dicono esattamente che la nostra estensione dei
numeri \`e consistente: poich\`e i numeri reali sono essenzialmente
definiti come limiti, il fatto che 
le operazioni sui limiti siano quelle
ovvie ci dice che i numeri reali si sommano, moltiplicano, dividono
etc. rispettando le propriet\`a gi\`a note per i numeri razionali.

\subsection{Esercizi}
\begin{enumerate}
\item Sia $a_0=1$, $a_1=1,3$, $a_2=1,33$, $a_3=1,333$, $a_4=1,3333$,
$a_5=1,33333$ e cos\`\i\ via. Si dimostri che 
\[
\lim_{n\to\infty}a_n=\frac 13.
\]
\item  Sia $a_0=0$, $a_1=0,9$, $a_2=0,99$, $a_3=0,999$, $a_4=0,9999$,
e cos\`\i\ via. Si dimostri che  
\[
\lim_{n\to\infty}a_n=1.
\]
\end{enumerate}
\end{document}